\chapter{Requirements and scope}
\label{ch:requirements}

% Budget and sources: docs/thesis/outline.md. Every claim about behaviour points at a test, a figure or a number in data/.

\section{Actors}
\label{sec:req-actors}

The system is offered as a service to institutions that handle disputes: banks that issue cards
and hold accounts, and payment providers acting for them. Each such institution is a
\emph{tenant}. A tenant's own systems, typically the online banking or card platform through which
a customer raises a complaint, open disputes and follow their progress through the API, and
identify themselves with a \emph{tenant key} issued for that purpose. The people who work the
disputes are the tenant's \emph{analysts}. They use the workbench, sign in through their own
institution's identity realm, belong to teams, and hold a rung of seniority within a team that
bounds what they may do. A \emph{tenant administrator} is an analyst who additionally manages
the tenant's keys, email templates and appearance. The \emph{customer}, finally, never uses the
system directly, but every deadline, letter and movement of money in it exists for them.

\section{Use cases}
\label{sec:req-use-cases}

One use case carries the system and the rest support it. A tenant system opens a dispute against a
transaction on the customer's account, giving the reason the customer chose. The regime is not
chosen by the caller; it follows from the transaction's rail and currency, so a card payment in
euro is handled under PSD2, a SEPA debit in euro under the direct-debit rules, a card payment in
dollars under Regulation~E and a credit-card payment in dollars under Regulation~Z
(\code{domain/rail.go}). The customer is told that the dispute has been opened. An analyst opens
the investigation, may send the customer a questionnaire and receive the answers, refunds or
credits the customer within the regime's clock, files a chargeback where the regime has one,
records its outcome, and closes the dispute with the settlement the outcome requires. A closed
dispute can be appealed once, which reopens the investigation with a new resolution clock.

Around that flow sit the use cases that make it workable for a team of people. Analysts find
disputes by identifier, state or reason, and can narrow the list to the ones whose deadlines have
passed. The customer's questionnaire depends on the dispute's reason, of which there are four: an
unauthorised payment, goods or services not received, a duplicate charge, and an amount different
from the one agreed. Analysts write to the customer from a catalogue of email templates, see every letter and email
sent, and can send one again. A team lead
can move a dispute to another team. Tenant administrators issue and revoke tenant keys and adapt the
email templates to their institution. Figure~\ref{fig:use-cases} traces each use case to the
operations of the API contract and to the section of chapter~\ref{ch:design} that designs it.

\begin{figure}[htbp]
\centering
\small
\begin{tabular}{L{4.6cm}L{2.6cm}L{4.4cm}l}
\toprule
Use case & Actor & Operations & Design \\
\midrule
Open a dispute & tenant system & \code{createDispute} & \ref{sec:design-state-machine} \\
Work a dispute through its lifecycle & analyst & \code{getDispute}, \code{applyDisputeEvent} & \ref{sec:design-state-machine} \\
Keep the regulatory deadlines & engine & (on every event) & \ref{sec:design-clocks} \\
Move money for a credit or reversal & engine & (on every event) & \ref{sec:design-ledger}, \ref{sec:design-core} \\
Collect the customer's account & analyst & \code{applyDisputeEvent} & \ref{sec:design-questionnaire} \\
Hold a credit on a risky dispute & engine, analyst & \code{applyDisputeEvent} & \ref{sec:design-fraud} \\
Find and filter disputes & analyst & \code{listDisputes} & \ref{sec:design-clocks} \\
Notify the customer & engine, analyst & \code{composeEmail}, \code{resendNotice} & \ref{sec:design-communications} \\
Move a dispute to another team & team lead & \code{reassignDispute} & \ref{sec:design-teams} \\
Manage keys and templates & tenant administrator & \code{createTenantKey}, \code{putTenantTemplate} & \ref{sec:design-tenancy} \\
\bottomrule
\end{tabular}
\caption{Use cases, the operations of the API contract that serve them (\code{docs/api/openapi.yaml}),
and the sections that design them.}
\label{fig:use-cases}
\end{figure}

\section{Quality requirements}
\label{sec:req-quality}

Four properties matter more than any single feature, because a failure of any of them is either a
breach of a legal duty or a loss of money. The first is correctness of the lifecycle: a dispute may
take no step its regime forbids, and every step the regime allows must be reachable and covered by
a test. The second is that money moves exactly once. A caller that repeats a request, because a
response was lost, must get the original answer rather than a second payment, and a refusal from
the core banking system must leave no trace of the step that caused it. The third is isolation
and accountability. No tenant may ever see another tenant's data, and this must hold even when
the application code is wrong, which means the database has to enforce it; the actor recorded on
every event must be the authenticated caller rather than whatever the request claims; and the
analyst who investigated a dispute may not also pay it out. The fourth is operability. Every
request must be traceable from the workbench to the database, and compliance with the deadlines
must be measurable as it happens, not reconstructed afterwards for a report.
Chapter~\ref{ch:evaluation} returns to each of the four with a test or a measurement.

\section{What was deliberately left out}
\label{sec:req-scope}

A dispute engine for a real bank would do more than this one, and the following were left out on
purpose. The rounds of a contested chargeback after the merchant's first answer, such as
pre-arbitration and arbitration, are folded into a single evidence step, since they add states
and deadlines but no new kind of decision. The extensions the regulations allow for complex
cases, such as Regulation~E's longer periods for new accounts, are not modelled; every dispute
runs the standard clocks. The card network and the core banking system are simulated behind
ports rather than integrated, as no real network or bank was available to a thesis project. Fraud
is scored with fixed, readable rules rather than a learned model, because there are no labelled
outcomes to learn from and a customer is entitled to an explanation. Each dispute has a single
currency and each tenant writes to its customers in a single language. PostgreSQL is the only
stateful dependency, with no separate cache or message broker, and there is no edge network in
front of the service, so the application does its own rate limiting. No real card numbers are
stored anywhere; all data is synthetic.

Two larger parts were marked from the outset as the first things to cut if time ran short: the
fraud-scoring engine and the analyst-facing frontend. Both survived, each in a deliberately modest
form, and chapters~\ref{ch:design} and~\ref{ch:implementation} explain what that form is and why.
